%=============================================================================
% WILSON RATIO AND ELECTROWEAK MIXING FORMULA CHAIN IN DFD
% Cross-reference trace: Where does the factor 6 come from?
% Generated: 2026-03-22
%=============================================================================

\documentclass[11pt]{article}
\usepackage{amsmath,amssymb,geometry,booktabs}
\geometry{margin=1in}
\title{The Wilson Ratio and Electroweak Mixing Chain in DFD:\\
Tracing the Factor of~6 in $\beta_{SU(2)}/\beta_{U(1)}$}
\author{Cross-Reference Audit}
\date{22 March 2026}
\begin{document}
\maketitle

%=============================================================================
\section{Executive Summary}
%=============================================================================

The factor 6 in the lattice coupling ratio
$\beta_{SU(2)}/\beta_{U(1)} = 6$
is \textbf{the product of two independently derived integers}:
\begin{equation}\label{eq:six}
\boxed{6 = \underbrace{\frac{n_2}{n_1}}_{=\,2}
      \;\times\;
      \underbrace{N_{\mathrm{gen}}}_{=\,3}}
\end{equation}

\begin{itemize}
\item $n_2/n_1 = 2$: the Frame Stiffness Theorem
      (Ricci curvature ratio on internal geometry).
\item $N_{\mathrm{gen}} = 3$: the index theorem on
      $\mathbb{C}P^2$ (Atiyah--Singer).
\end{itemize}

It is \emph{not} a lattice convention. It is \emph{not} the Wilson
normalization factor of 4 (that enters separately in $g_2^2$).
It is not $2+4$.  The factorisation is $2 \times 3$, period.

%=============================================================================
\section{The Complete Formula Chain}
%=============================================================================

\subsection{Step 1: Frame Stiffness Theorem (Appendix~F, Theorem~F.13)}

Gauge fields arise as Berry connections on the internal manifold
$M_{\mathrm{int}} = \mathbb{C}P^2 \times S^3$.
The gauge sectors correspond to isometries acting on complex subspaces
$V_r$ of dimension $n_r$ within the $(3,2,1)$ partition.

The energy cost of a unit frame rotation scales with the Ricci curvature
of the corresponding projective space $\mathbb{C}P^{n_r-1}$:
\[
  R_{i\bar{j}}(\mathbb{C}P^{n_r-1}) = n_r \cdot g_{i\bar{j}}^{\mathrm{FS}},
\]
giving the frame stiffness law:
\begin{equation}\label{eq:kappa}
\kappa_r = n_r\,\kappa_0,
\end{equation}
where $n_{U(1)}=1$, $n_{SU(2)}=2$, $n_{SU(3)}=3$.

Therefore:
\begin{equation}\label{eq:stiffness-ratio}
\boxed{\frac{\kappa_{U(1)}}{\kappa_{SU(2)}} = \frac{n_1}{n_2}
= \frac{1}{2}}
\end{equation}

\textbf{Source files:}
\begin{itemize}
\item \texttt{appendix\_F.tex}, Theorem~F.13 (``Frame Stiffness from Geometry''),
      lines 355--383.
\item \texttt{section\_quantum.tex}, \S13.3 (``Yang--Mills Kinetic Terms from
      Frame Stiffness'').
\item PDF paper, Eq.~(8)--(9), p.~7.
\end{itemize}

Lattice verification: $\kappa_{U(1)}/\kappa_{SU(2)} = 0.495 \pm 0.020$
(PDF, Section~V\,E, p.~17).

\subsection{Step 2: Generation Count from Index Theorem}

The number of chiral fermion generations is determined by the
flux-product rule on $\mathbb{C}P^2 \times S^3$
(Appendix~F, Theorem~F.8):
\[
  N_{\mathrm{gen}} = |k_3 \cdot k_2 \cdot q_1|,
\]
where $(k_3,k_2,q_1)=(1,1,3)$ is the unique energy-minimising flux
configuration satisfying the Spin$^c$ constraint $q_1=3$.
Therefore:
\begin{equation}\label{eq:Ngen}
\boxed{N_{\mathrm{gen}} = 3.}
\end{equation}

\textbf{Source files:}
\begin{itemize}
\item \texttt{appendix\_F.tex}, Lemma~F.4 ($q_1=3$), Theorem~F.7
      (energy minimisation), Theorem~F.8 (flux-product rule).
\item \texttt{section\_quantum.tex}, \S13.5 (``Generation Counting'').
\item PDF paper, abstract item~3 and Eq.~(17), p.~10.
\end{itemize}


\subsection{Step 3: The Wilson Ratio = Stiffness Ratio $\times$ Generation Count}

The lattice action ratio $\beta_{SU(2)}/\beta_{U(1)}$ is set by two
independent DFD quantities:
\begin{equation}\label{eq:wilson-ratio}
\boxed{\frac{\beta_{SU(2)}}{\beta_{U(1)}}
  = \frac{n_2}{n_1} \times N_{\mathrm{gen}}
  = 2 \times 3 = 6.}
\end{equation}

\textbf{Why $N_{\mathrm{gen}}$ enters:} All three fermion generations
contribute equally to the effective lattice coupling because the
Atiyah--Singer index on $\mathbb{C}P^2$ forces exactly three chiral
families. Each generation adds one unit of the SU(2) representation
content to the effective plaquette action, multiplying the bare
stiffness ratio by 3.

\textbf{Source files:}
\begin{itemize}
\item \texttt{appendix\_K.tex}, \S K.2 ``(4) Wilson ratio from topology'',
      lines 130--135.
\item PDF paper, abstract item~3; Eq.~(17), p.~10; Section~III\,C
      ``Constraint~3: The lattice ratio from topology'', pp.~10--11.
\end{itemize}

Lattice verification: ten ratios tested (3, 4, 5, 5.5, \textbf{6}, 6.25,
6.5, 7, 8, 9). Only ratio 6 yields $\alpha = 1/137$. Fractional values
5.5, 6.25, 6.5 also fail, proving the factor must be \emph{exactly} 6.
(\texttt{appendix\_K.tex}, Table~K.3; PDF, Table~III, p.~11.)


\subsection{Step 4: Wilson Normalization of Gauge Couplings}

With the standard Wilson action normalization
$\mathrm{Tr}(T^a T^b) = \frac{1}{2}\delta^{ab}$ for SU(2),
the Wilson plaquette action gives $\beta_{SU(2)} = 4/g_2^2$.
For compact U(1), $\beta_{U(1)} = 1/g_1^2$.

The stiffnesses $\kappa_r$ (measured non-perturbatively via the
background-field method) relate to the gauge couplings through:
\begin{equation}\label{eq:g-from-kappa}
g_1^2 = \frac{1}{\kappa_{U(1)}}, \qquad
g_2^2 = \frac{4}{\kappa_{SU(2)}}.
\end{equation}

The factor of 4 in $g_2^2 = 4/\kappa_{SU(2)}$ is the standard Wilson
action normalization for SU(2), arising from the trace convention
$\mathrm{Tr}(T^a T^b) = \tfrac{1}{2}\delta^{ab}$.
This is \emph{separate} from the ratio-of-6 discussed above.

\textbf{Source files:}
\begin{itemize}
\item PDF paper, Eq.~(12), p.~8; Listing~2, p.~24.
\item \texttt{appendix\_Z\_complete\_params.tex}, lines 34--37.
\item \texttt{supplementary/PRL\_SM\_from\_Topology.tex}, lines 70--74.
\end{itemize}


\subsection{Step 5: Electroweak Mixing and $\alpha$ Extraction}

Electromagnetism emerges from mixing of U(1) and neutral SU(2):
\begin{equation}\label{eq:em-mixing}
\frac{1}{e^2} = \frac{1}{g_1^2} + \frac{1}{g_2^2}
= \kappa_{U(1)} + \frac{\kappa_{SU(2)}}{4}.
\end{equation}

Substituting $\kappa_{SU(2)} = 2\kappa_{U(1)}$ (Frame Stiffness Theorem):
\begin{equation}
\frac{1}{e^2} = \kappa_{U(1)} + \frac{2\kappa_{U(1)}}{4}
= \kappa_{U(1)} + \frac{\kappa_{U(1)}}{2}
= \frac{3}{2}\,\kappa_{U(1)}.
\end{equation}

The fine-structure constant is then:
\begin{equation}\label{eq:alpha-from-kappa}
\alpha_W = \frac{e^2}{4\pi}
= \frac{1}{4\pi \cdot \frac{3}{2}\kappa_{U(1)}}
= \frac{1}{6\pi\,\kappa_{U(1)}},
\end{equation}
i.e.,
\begin{equation}\label{eq:alpha-inverse}
\boxed{\alpha^{-1} = 6\pi\,\kappa_{U(1)}.}
\end{equation}

\textbf{Source files:}
\begin{itemize}
\item PDF paper, Eq.~(10), (12), (13), pp.~7--8; Listing~2, p.~24.
\end{itemize}


%=============================================================================
\section{Dissecting the Factor of 6 in $\alpha^{-1} = 6\pi\kappa$}
%=============================================================================

The ``6'' in $\alpha^{-1} = 6\pi\,\kappa_{U(1)}$ has a completely
transparent algebraic origin.  It is:
\begin{equation}
6 = \underbrace{4}_{\text{Wilson norm}}
    \times \underbrace{\frac{1}{\kappa_{U(1)}/\kappa_{SU(2)}}}_{\text{inverse
    stiffness ratio}}
    \times \frac{1}{4}
    + 1
    \quad\longleftarrow\;\text{NO.  This is wrong.}
\end{equation}

The correct decomposition is elementary algebra from
$1/e^2 = 1/g_1^2 + 1/g_2^2$:
\begin{align}
\frac{1}{e^2}
&= \kappa_{U(1)} + \frac{\kappa_{SU(2)}}{4} \nonumber\\
&= \kappa_{U(1)} + \frac{2\kappa_{U(1)}}{4} \nonumber\\
&= \kappa_{U(1)}\!\left(1 + \frac{1}{2}\right) \nonumber\\
&= \frac{3}{2}\,\kappa_{U(1)}.
\end{align}

So:
\[
\alpha^{-1} = 4\pi/e^2 = 4\pi \cdot \frac{3}{2}\kappa_{U(1)}
= 6\pi\,\kappa_{U(1)}.
\]

\textbf{The 6 in $\alpha^{-1} = 6\pi\kappa$ is therefore
$4\pi \times 3/2 = 6\pi$, where:}
\begin{itemize}
\item The $4\pi$ is the definition $\alpha = e^2/(4\pi)$.
\item The $3/2$ is $1 + n_2/(2\cdot 2) = 1 + 1/2$, coming from
      $1/e^2 = \kappa_{U(1)} + \kappa_{SU(2)}/4$ with the stiffness
      ratio $\kappa_{SU(2)}/\kappa_{U(1)} = 2$.
\end{itemize}

Stated differently, the coefficient $3/2$ in $1/e^2 = (3/2)\kappa_{U(1)}$
decomposes as:
\begin{equation}
\frac{3}{2} = \underbrace{1}_{\substack{\text{U(1) contribution}\\
  1/g_1^2 = \kappa_{U(1)}}}
+ \underbrace{\frac{1}{2}}_{\substack{\text{SU(2) contribution}\\
  1/g_2^2 = \kappa_{SU(2)}/4 = \kappa_{U(1)}/2}}
\end{equation}

The $1/2$ from SU(2) arises because:
\begin{itemize}
\item $\kappa_{SU(2)} = 2\kappa_{U(1)}$ (stiffness ratio from Ricci
      curvature, i.e., $n_2/n_1 = 2$).
\item But $1/g_2^2 = \kappa_{SU(2)}/4$ (Wilson normalization, from
      $\mathrm{Tr}(T^aT^b) = \tfrac12\delta^{ab}$).
\item Net: $1/g_2^2 = 2\kappa_{U(1)}/4 = \kappa_{U(1)}/2$.
\end{itemize}

\textbf{This is NOT the same 6 as in the Wilson ratio
$\beta_{SU(2)}/\beta_{U(1)} = 6$.}
The coincidence that both equal 6 is accidental from the standpoint of
the algebra.  The Wilson ratio 6 = $2 \times 3$ (stiffness ratio
$\times$ generation count).  The coefficient in
$\alpha^{-1} = 6\pi\kappa$ comes from $4\pi \times 3/2$
(definition of $\alpha$ $\times$ electroweak mixing coefficient).


%=============================================================================
\section{The Two Distinct ``Factor of 6'' Appearances}
%=============================================================================

\begin{center}
\renewcommand{\arraystretch}{1.4}
\begin{tabular}{lll}
\toprule
\textbf{Quantity} & \textbf{Value} & \textbf{Origin} \\
\midrule
$\beta_{SU(2)}/\beta_{U(1)}$ & $6 = 2 \times 3$
  & Stiffness ratio $n_2/n_1$ $\times$ generation count $N_{\mathrm{gen}}$ \\
$\alpha^{-1}/(\pi\kappa_{U(1)})$ & $6 = 4\pi \times 3/(2\pi \times 2)$
  & Definition of $\alpha$ $\times$ electroweak mixing coefficient \\
\bottomrule
\end{tabular}
\end{center}

The first 6 sets the \emph{lattice input parameters}
$(\beta_{U(1)}, \beta_{SU(2)}) = (3.80, 22.80)$.

The second 6 converts the \emph{measured output} $\kappa_{U(1)}$
into $\alpha$.

They are logically independent.  The lattice simulation bridges between
them: it takes the input betas (set by the Wilson ratio of 6) and
produces output kappas, from which $\alpha$ is extracted via the
electroweak mixing formula (which happens to also involve a factor 6).

%=============================================================================
\section{The Full Derivation Chain (Summary)}
%=============================================================================

\begin{equation}
\boxed{
\begin{aligned}
&\text{Topology:}\quad k_{\max} = \chi(\mathbb{C}P^2, \mathcal{O}(9) \oplus \mathcal{O}^{\oplus 5}) = 60 \\
&\text{CS vacuum:}\quad \beta_{U(1)} = \langle k+2 \rangle_{k_{\max}=60} = 3.80 \\
&\text{Frame Stiffness:}\quad \kappa_{U(1)}/\kappa_{SU(2)} = n_1/n_2 = 1/2 \\
&\text{Index Theorem:}\quad N_{\mathrm{gen}} = 3 \\
&\text{Wilson ratio:}\quad \beta_{SU(2)}/\beta_{U(1)} = (n_2/n_1) \times N_{\mathrm{gen}} = 2 \times 3 = 6 \\
&\Rightarrow\; \beta_{SU(2)} = 6 \times 3.80 = 22.80 \\
&\text{Lattice MC at } (\beta_{U(1)},\beta_{SU(2)}) = (3.80, 22.80)
  \;\Rightarrow\; \kappa_{U(1)} \approx 7.25 \\
&\text{Wilson norms:}\quad g_1^2 = 1/\kappa_{U(1)},\quad g_2^2 = 4/\kappa_{SU(2)} \\
&\text{EM mixing:}\quad 1/e^2 = 1/g_1^2 + 1/g_2^2 = (3/2)\,\kappa_{U(1)} \\
&\text{Result:}\quad \alpha^{-1} = 4\pi/e^2 = 6\pi\,\kappa_{U(1)} \approx 137
\end{aligned}
}
\end{equation}

%=============================================================================
\section{Where Each Factor Comes From}
%=============================================================================

\begin{center}
\renewcommand{\arraystretch}{1.3}
\begin{tabular}{llll}
\toprule
\textbf{Factor} & \textbf{Value} & \textbf{Source} & \textbf{File} \\
\midrule
$n_1$ & 1 & dim $V_{U(1)} = \mathbb{C}P^0$ & appendix\_F.tex \\
$n_2$ & 2 & dim $V_{SU(2)} = \mathbb{C}P^1$ & appendix\_F.tex \\
$n_2/n_1$ & 2 & Frame Stiffness Thm (Ricci curv.) & appendix\_F.tex, Thm F.13 \\
$N_{\mathrm{gen}}$ & 3 & $|k_3 k_2 q_1| = |1 \cdot 1 \cdot 3|$ & appendix\_F.tex, Thm F.8 \\
$\beta_{SU(2)}/\beta_{U(1)}$ & 6 & $(n_2/n_1) \times N_{\mathrm{gen}}$ & appendix\_K.tex, line 133 \\
Wilson norm (SU(2)) & 4 & $\mathrm{Tr}(T^aT^b) = \frac{1}{2}\delta^{ab}$ & appendix\_Z.tex, line 34 \\
$3/2$ in $1/e^2$ & $3/2$ & $1 + \kappa_{SU(2)}/(4\kappa_{U(1)}) = 1 + 1/2$ & PDF Eq.~(13) \\
$4\pi$ & $4\pi$ & definition $\alpha = e^2/(4\pi)$ & standard QFT \\
$6\pi$ in $\alpha^{-1}$ & $6\pi$ & $= 4\pi \times 3/2$ & PDF Listing~2 \\
\bottomrule
\end{tabular}
\end{center}

%=============================================================================
\section{The Weinberg Angle from the Same Geometry}
%=============================================================================

The Weinberg angle also follows from the partition stiffnesses.
With canonical normalizations (Appendix~Z, Theorem~Z.1):
\[
g'^{-2} \propto \kappa_1 \cdot \mathrm{Tr}(Y^2) = \kappa_1 \cdot \frac{10}{3},
\qquad
g^{-2} \propto \kappa_2 \cdot \frac{1}{2},
\]
giving:
\[
\frac{g'^2}{g^2} = \frac{\kappa_2 \cdot (1/2)}{\kappa_1 \cdot (10/3)}
= \frac{2 \cdot (1/2)}{1 \cdot (10/3)} = \frac{3}{10},
\]
and therefore:
\begin{equation}
\sin^2\theta_W = \frac{g'^2}{g^2 + g'^2} = \frac{3/10}{1 + 3/10}
= \frac{3}{13} = 0.2308.
\end{equation}

Measured: $\sin^2\theta_W(\overline{\mathrm{MS}}, M_Z) = 0.23122$
(agreement: 0.2\%).

\textbf{Source:} \texttt{appendix\_Z\_complete\_params.tex}, lines 18--54;
\texttt{supplementary/PRL\_SM\_from\_Topology.tex}, lines 50--93.

%=============================================================================
\section{Conclusions}
%=============================================================================

\begin{enumerate}
\item The Wilson ratio $\beta_{SU(2)}/\beta_{U(1)} = 6$ is the product
      of two topological integers: the stiffness ratio $n_2/n_1 = 2$
      (from Ricci curvature on the internal manifold) and the generation
      count $N_{\mathrm{gen}} = 3$ (from the index theorem on
      $\mathbb{C}P^2$).

\item The coefficient 6 in $\alpha^{-1} = 6\pi\kappa_{U(1)}$ is
      $4\pi \times 3/2$, where the $3/2$ comes from electroweak mixing
      with the stiffness ratio $\kappa_{SU(2)} = 2\kappa_{U(1)}$ and
      the Wilson normalization factor of 4 for SU(2).

\item These are two \emph{different} appearances of the number 6.
      Their numerical coincidence is algebraically traceable but not
      a single shared origin.

\item The Wilson normalization factor of 4 (from
      $\mathrm{Tr}(T^aT^b) = \frac{1}{2}\delta^{ab}$) enters
      $g_2^2 = 4/\kappa_{SU(2)}$ but does \emph{not} enter the Wilson
      ratio of 6.  These are independent.

\item Ten lattice ratios were tested (including fractional values).
      Only the integer 6 works.  This is strong evidence that both
      constituent factors (2 and 3) are exactly correct.
\end{enumerate}

\end{document}
